\section{Discussion}

\subsection{Answer to the research question}
In this environment, the time of the basic $O(n^3)$ multiplication grows close to cubically for small and medium
matrices ($k\approx3.2$ for $50\le n\le400$ in all three languages), but the cost per multiply--add \emph{increases}
for larger matrices in the compiled/JIT languages: in C from $1.0$--$1.4$\,ns ($200\le n\le1000$, excluding $n=512$) to $6.8$\,ns
($n=3000$), in Java from $\approx1.6$\,ns to $11.1$\,ns ($n=2048$); in Python it rises from $50$--$60$\,ns to
$\approx100$\,ns. Python is $49$--$75$ times slower than C and Java at $50\le n\le1000$. C and Java are
comparable for $50\le n\lesssim1000$ (ratio $0.77$--$1.51$), after which the \texttt{double[n][n]} Java program falls behind
($2.2$--$4.3\times$). Memory is explained by array storage plus a runtime baseline: $24$, $\approx22$ (Java, $n\le1500$) and
$\approx121$ bytes per $n^2$, with baselines of $1.8$, $45.7$ and $9.8$\,MiB for C, Java and Python. Time, not memory, limits
the feasible size: $n=3000$ (C), $2048$ (Java), $1250$ (Python) under a $300$\,s-per-multiplication limit.

\subsection{Hypotheses}
\begin{itemize}
  \item \textbf{H1 (cubic growth)}: supported only partially. The exponent is $3.17$--$3.27$ for $n\le400$, but $3.29$ (Python),
        $3.63$ (C) and $4.01$ (Java) over the larger windows. The operation count is cubic; the \emph{time} is not, because the cost
        per operation depends on $n$ (Figure~\ref{fig:norm}).
  \item \textbf{H2 (C $\lesssim$ Java $\ll$ Python)}: Python is $49$--$75\times$ slower for $50\le n\le1000$, consistent with ``one to two orders
        of magnitude''. C $\lesssim$ Java holds only for $50\le n\lesssim1000$ (within noise); beyond that Java is up to $4.3\times$ slower.
  \item \textbf{H3 (cost per operation rises with $n$)}: supported in all three languages (C $0.5\to6.8$\,ns, Java $0.9\to11.1$\,ns,
        Python $50\to106$\,ns), with additional spikes at powers of two.
  \item \textbf{H4 (memory model)}: supported for C ($24.0$ B$/n^2$) and Python ($120.9$ B$/n^2$); Java matches ($22.3$) up to $n=1500$ but
        needs $\approx32$ B$/n^2$ at $n\ge2000$.
\end{itemize}

\subsection{Interpretation, and how strong the evidence is}
\paragraph{Python.} The $\approx50$--$60$\,ns per multiply--add for $n\le200$ is the cost of executing the inner loop in the
interpreter, where each iteration performs several bytecode dispatches, two double indexings and the allocation of a float object;
this is consistent with the observed constant per-operation cost at small $n$ and the $\approx50$--$75\times$ gap, but the cost
of individual bytecodes was not measured here. The further increase to $82$--$106$\,ns at $n\ge400$ coincides with the working set
growing beyond the caches (about $40$ bytes per element, Table~\ref{tab:mem}), where every access dereferences a pointer to a float
object; Python's increase ($\approx2\times$) is smaller than in C ($\approx5\times$) because interpreter overhead is a large constant part of its cost.

\paragraph{C: powers of two and large $n$.} Two effects are visible. (i) At $n=512$ and $1024$ the matrix row stride is a multiple of
$4096$ bytes, so the successive elements $b_{kj}$ of a column are mapped to few cache sets; the time per multiply--add jumps by
$\approx1.8$--$1.9\times$ relative to $n\pm$ a few (Figure~\ref{fig:pow2}). The effect appears in C and in Java with a flat array,
and not in Java with \texttt{double[n][n]}, which is consistent with a memory-layout explanation (row addresses of separately allocated
Java arrays are not exact multiples of the page size) \cite{drepper,hennessy}. (ii) For $n\gtrsim1250$ the cost per operation grows
steadily (Figure~\ref{fig:norm}). The matrix $B$ occupies $8n^2$ bytes ($8$\,MB at $n=1000$, $32$\,MB at $n=2048$, $72$\,MB at $n=3000$),
which exceeds the L2 cache ($2$\,MiB per core) and approaches or exceeds the L3 cache ($33$\,MiB, shared) of this machine, and a column
walk touches a different $4$\,KiB page at every step for $n\ge512$, which also stresses the address-translation caches. These are plausible
causes, consistent with the trend, but \textbf{they were not measured}: no hardware performance counters were available in the virtual
machine, so cache misses and TLB misses were not observed.

\paragraph{Java versus C.} The gap in Figure~\ref{fig:time} is not a property of the JIT arithmetic: Java with a flat array is within
$1.07$--$1.49\times$ of C for $500\le n\le2000$, whereas the \texttt{double[n][n]} version is up to $4.3\times$ slower, and the
jagged/flat ratio grows with $n$. A jagged array adds, for every $b[k][j]$, a load of the row reference and a length check on a
\emph{separate} object, and rows are separate memory blocks; at large $n$ this plausibly increases cache and TLB pressure. The experiment
supports ``layout explains most of the gap'' but does not identify the mechanism, and the remaining $7$--$49\,\%$ (bounds checks,
JIT code quality, or other effects) was not isolated. Warm-up played a minor role: the first warm-up sample was at most $1.6\times$ the steady
value and the measured samples were stable (Figure~\ref{fig:var}).

\paragraph{C compiler options.} \texttt{-O0} is $2.2$--$3.7\times$ slower; \texttt{-O1} to \texttt{-O3} are indistinguishable within the noise.
Because floating-point reductions cannot be re-associated without \texttt{-ffast-math} and the column access of $B$ is strided, the
compiler has little to improve; \texttt{-march=native} did not help either. This supports using \texttt{-O2} as a representative optimised
configuration for this particular kernel, nothing more.

\subsection{Limitations}
\begin{itemize}
  \item \textbf{One virtual machine.} All data come from one $2$-vCPU KVM guest on a shared host: other tenants, frequency scaling and
        host scheduling are not observable or controlled, and the processes were not pinned. Absolute times are not transferable to other machines.
  \item \textbf{Between-process variability.} At $n=1024$ the median of six independent processes varies by $13\,\%$ (C) and $40\,\%$ (Java)
        and the within-process spread of Java is also large. Most comparisons of ratios close to $1$ at a single size (e.g.\ Java/C
        $=0.77$ at $n=512$ or $0.96$ at $n=1024$) are therefore \emph{not} significant. The IQRs in Table~\ref{tab:time} describe
        within-process repetitions only.
  \item \textbf{Few repetitions at large $n$} (2--3, forced by the budget), a single pair of input matrices per size, and a single run
        order per size (interleaved by language to reduce, not eliminate, time-of-day effects).
  \item \textbf{Exploratory supplementary experiments.} They were designed after seeing the main results, have one process per cell (except $n=1024$) and
        test explanations indirectly; the mechanisms above remain hypotheses.
  \item \textbf{Scope.} Results apply to gcc $13.3$ \texttt{-O2}, OpenJDK $21$ with the serial GC, CPython $3.13$ and the algorithm, layouts and
        inputs described here. They do not support a statement that one language is faster than another in general; other compilers, JVM settings, Python
        implementations or matrix libraries could change the picture drastically (and are the topic of the next assignments).
  \item \textbf{Memory metric.} RSS is not the same as heap use or allocated storage; the Java excess at $n\ge2000$ was not investigated.
\end{itemize}

\section{Conclusions}
\begin{enumerate}
  \item Under the described conditions, the basic triple-loop multiplication takes $\approx n^{3.2}$ for $n\le400$ in C, Java and
        Python, and the cost per multiply--add rises for larger matrices (C up to $6.8$\,ns at $n=3000$, Java $11.1$\,ns at
        $n=2048$, Python $\approx100$\,ns at $n\approx1000$--$1250$).
  \item Python is $49$--$75\times$ slower than C and Java for $50\le n\le1000$. C and Java agree within a factor of $1.5$ for $50\le n\lesssim1000$;
        beyond, \texttt{double[n][n]} Java is $2.2$--$4.3\times$ slower than C, and a flat-array Java variant removes most of that gap ($\le1.5\times$ of C).
  \item C (and flat Java) is $\approx1.7$--$1.9\times$ slower at $n=512$ and $n=1024$ than at neighbouring sizes.
  \item Memory follows array storage plus a runtime baseline: $24$ B$/n^2$ (C), $\approx22$ (Java, $n\le1500$), $\approx121$ (Python), with baselines $1.8$, $45.7$ and $9.8$\,MiB.
  \item With a $300$\,s per-multiplication limit the largest feasible sizes were $3000$ (C), $2048$ (Java) and $1250$ (Python); the first
        failures ($4096$, $3000$, $1500$) were timeouts, never memory failures.
  \item Unresolved: the hardware mechanisms behind the large-$n$ slowdown and the power-of-two spikes (needs performance counters),
        the origin of Java's extra memory at $n\ge2000$, and the effect of the machine (needs repetition on other hardware).
\end{enumerate}
The numbers apply to this experiment only; the optimisation techniques that address several of the observed effects (loop order, blocking, sparse
formats) are left for Assignment~2.

\section{Use of generative AI}
In line with the course policy, the use of AI is disclosed: the tool was \textbf{Claude (Anthropic)}. It was used for planning the experimental
design, writing the benchmark driver, the test and analysis scripts, the supplementary experiments and a first draft of this report.
The affected parts are \texttt{scripts/}, \texttt{src/} and \texttt{report/}. The author is responsible for understanding the code, validating the measurements,
checking the references and defending the conclusions; generated text is not evidence, and every number in this report was produced by the included code
from the raw data in \texttt{results/}.

\section{Reproducibility materials}
\label{sec:materials}
Source code, build and run instructions, raw measurements, input-generation scripts and analysis scripts are in the repository
\repourl{} and in \path{Assignment01_GonzalezFernandez_Materials.zip}. The \texttt{README.md} gives the exact commands
(\texttt{make build inputs test bench cflags analyze report}); the input matrices are not stored but regenerated with
\texttt{scripts/gen\_inputs.py} and verified against the SHA-256 digests in \texttt{data/manifest.json}. Raw data:
\path{results/raw_measurements.csv} (one row per sample, warm-up included) and \path{results/runs.csv} (one row per process);
environment: \path{results/environment.json}; every figure and table is generated by \path{scripts/analyze.py}.
The basic triple-loop structure was taken from the example implementations (Python, Java, Rust and C, $1024\times1024$) that were given to the author as examples for this
assignment; their timings were not used. The shared binary inputs, the scalar-accumulator kernels, monotonic timers, measurement driver, validation, supplementary experiments and analysis were
written for this study with the AI assistance disclosed above. No other external code, text or figures are reused. All numbers in this report come from the run documented in
\texttt{results/}.

\begin{thebibliography}{99}
\bibitem{clrs} T.~H.~Cormen, C.~E.~Leiserson, R.~L.~Rivest and C.~Stein, \emph{Introduction to Algorithms}, 4th~ed. Cambridge, MA, USA: MIT Press, 2022.
\bibitem{drepper} U.~Drepper, ``What every programmer should know about memory,'' Red~Hat, Inc., Nov. 2007.
\bibitem{hennessy} J.~L.~Hennessy and D.~A.~Patterson, \emph{Computer Architecture: A Quantitative Approach}, 6th~ed. Cambridge, MA, USA: Morgan Kaufmann, 2017.
\bibitem{higham} N.~J.~Higham, \emph{Accuracy and Stability of Numerical Algorithms}, 2nd~ed. Philadelphia, PA, USA: SIAM, 2002.
\bibitem{goldberg} D.~Goldberg, ``What every computer scientist should know about floating-point arithmetic,'' \emph{ACM Comput. Surv.}, vol.~23, no.~1, pp.~5--48, Mar. 1991.
\bibitem{georges} A.~Georges, D.~Buytaert and L.~Eeckhout, ``Statistically rigorous Java performance evaluation,'' in \emph{Proc. 22nd ACM SIGPLAN Conf. Object-Oriented Programming Systems, Languages and Applications (OOPSLA)}, Montreal, Canada, 2007, pp.~57--76.
\bibitem{hoefler} T.~Hoefler and R.~Belli, ``Scientific benchmarking of parallel computing systems: Twelve ways to tell the masses when reporting performance results,'' in \emph{Proc. Int. Conf. High Performance Computing, Networking, Storage and Analysis (SC '15)}, Austin, TX, USA, 2015, Art.~no.~73.
\bibitem{numpy} C.~R.~Harris \emph{et al.}, ``Array programming with NumPy,'' \emph{Nature}, vol.~585, pp.~357--362, 2020, doi: 10.1038/s41586-020-2649-2.
\bibitem{py312} Python Software Foundation, ``What's new in Python 3.12,'' Python documentation (built-in \texttt{sum()} uses compensated summation for floats; CPython issue gh-100425).
\end{thebibliography}
